\part{自救方法论：顺着大脑参数设计生活}

\section{认知自救：先让自己舒服}

进入方法之前，必须先打地基。因为一个残酷的事实是：\term{方法再好，一边做一边骂自己「我怎么这么懒」，方法也会失效}。

ADHD 自救的顺序，永远是认知在前，操作在后。这一章先把「舒服」这件事说透。

\subsection{从「我有病」到「我有另一套参数」}

你大概在无数个深夜问过自己：我是不是太懒？我是不是没救了？我是不是根本不适合活着这件事？

把这些问题翻译一下，它们都建立在同一个前提上：\term{「我应该像正常人那样运转，而我坏了」}。这个前提一旦松动，整座山都跟着松。

换个说法试试：你不是坏了，你是\term{参数不同}。参数没有对错，只有适配。GPU 跑不了单线程是「缺陷」吗？不是，是它被用错了地方。你过去所有「失败」，都是这套参数在错误的运行环境里的正常结果——不是人品问题，不是能力问题，是\term{不匹配}。

\begin{insight}
病耻感是最贵的内耗。它让你把大量的执行力花在「证明自己没坏」上，而不是花在「把自己用好」上。放下「我有病」三个字，是自救的第一个动作。
\end{insight}

\subsection{停止道德账本}

ADHD 群体有一种特有的内耗：随身携带一本道德账本，每拖延一次，就记上一笔「我又懒了」「我又辜负了期待」。到晚上临睡前，把账本翻出来，一边数一边骂自己。

问题在于：\term{自责不是动力，是阻力}。你花在自责上的精力，够你做三件事了。更糟的是，自责会推高皮质醇——还记得第三章吗？皮质醇一高，启动就锁死，于是你陷入「拖—骂—更拖」的死循环。

\begin{checklist}
  \item 把「我怎么这么懒」改成「今天的燃料又不对路」——前者是审判，后者是观察；
  \item 把「我必须今天完成」改成「我今天先做五分钟」——前者是警报，后者是入口；
  \item 每天晚上，把道德账本换成「今天大脑运行日志」：什么时段状态好、什么刺激管用、什么环境有害——这是数据，不是罪证。
\end{checklist}

\subsection{自我关怀不是放纵}

很多人不敢「放过自己」，是因为怕一放就彻底躺平。这里要分清两件事：

\term{放纵}是「今天算了」，是关闭；\term{自我关怀}是「今天这样也是合理的，明天换个办法」，是重新开始。

差别在哪？放纵之后更恨自己；自我关怀之后更有力气。羞耻会锁死执行功能，接纳会解锁。对 ADHD 来说，自我关怀不是奖赏，是\term{生产力}——它是让大脑敢再试一次的前提。

\subsection{接受「我需要外部脚手架」}

眼镜不可耻，外骨骼也不可耻。

别人靠内建自控（前额叶直接指挥），你靠外置结构（环境、清单、提示、陪伴）。只要结果一样——作业交了、事情做了、生活没塌——\term{谁在乎结构装在大脑里还是装在环境里}？

很多学历上层的 ADHD 在这里卡住：太要强，觉得「靠环境布置来让自己做事」是作弊，显得自己不行。恰恰相反，\term{懂得布置环境，才是对这套算法最深的尊重}。把最重要的文件放桌面最显眼处、把手机放另一个房间、让记录一直开着——这些不是逃避，是工程。

\subsection{允许「蒙混过关」和「放弃」是合法选项}

这是原视频里最有争议、也最被低估的一招（BV1PGcZzsE9K，BV1Y2wizyEua）：对 ADHD 来说，\term{「蒙混过关」和「放弃」应该是合法选项}。

先说放弃：有些事你实在拖太久了，而且它对生活没有绝对影响，那就干脆决定不做。原视频的观察是：\term{你决定不做了，之后重新产生冲动的概率反而更高}——因为「必须做」的警报撤了，大脑不再对抗，兴趣才有机会重新长出来。先混过去，再放下，再重新自发产生兴趣，这个循环比被任务死死压住健康得多。

再说蒙混过关：接到一个外界派下来的线性任务，别从「一丝不苟地做下来」开始（那对 ADHD 是死刑），先想「怎么把它蒙混过关」——这个过程反而会产出很多想法，再迭代成能交的版本。注意，这招涉及社交羞耻（「我一开始就在想蒙混过关，会不会显得我很坏」），心理负担极大，没几个人愿意用，但它确实是存在的、有效的选项。

\begin{notehot}
这两招的真正用途，不是让你变油滑，而是\term{给你腾出重新启动的空间}。对一个启动困难的大脑，「允许不做」和「允许烂做」是两把钥匙：钥匙插进去，门才开得动。等重新有了兴趣和能量，再好好做，那时候才是真质量。
\end{notehot}

\subsection{一个微妙的平衡}

当然也要说清楚：以上都不是「给自己的借口」——借口是逃避之后的自我安慰，这些是\term{给启动让路的工程手段}。区别只有一个检验标准：\term{用完之后，你是更无力了，还是更有力气了}。更无力，那是借口；更有力气，那是方法。请用这个标准，而不是道德标准，来判断自己。

\key{舒服不是躺平，舒服是停止和重力对抗。认知上彻底放下「我应该像正常人一样」，方法才能落地。}
